\documentclass[11pt,letterpaper]{article}
\usepackage[T1]{fontenc}
\usepackage[utf8]{inputenc}
\usepackage[margin=0.91in,headheight=15pt,headsep=18pt,footskip=27pt]{geometry}
\usepackage{newpxtext}
\usepackage{amsmath,amsthm}
\usepackage{newpxmath}
\usepackage{mathtools,microtype}
\usepackage{booktabs,tabularx,longtable,array}
\usepackage{enumitem}
\usepackage{xcolor,titlesec,fancyhdr}
\usepackage{xurl}
\definecolor{Navy}{RGB}{25,49,72}
\definecolor{Steel}{RGB}{47,81,111}
\definecolor{Muted}{RGB}{90,99,107}
\usepackage[colorlinks=true,linkcolor=Steel,citecolor=Steel,urlcolor=Steel,
  bookmarksnumbered=true,pdfencoding=auto,psdextra]{hyperref}
\hypersetup{pdftitle={Finite Realization of Coarse-Information Profiles},
 pdfauthor={Research draft prepared for Vladimir Reshetnikov with ChatGPT},
 pdfsubject={Countable Turing ideals, coarse descriptions, generic secret sharing, and representative spectra}}
\urlstyle{same}
\setlength{\parindent}{1.2em}
\setlength{\parskip}{3pt plus .5pt minus .4pt}
\setlength{\emergencystretch}{2em}
\linespread{1.035}
\titleformat{\section}{\Large\bfseries\color{Navy}}{\thesection}{.65em}{}
\titleformat{\subsection}{\large\bfseries\color{Navy}}{\thesubsection}{.6em}{}
\titleformat{\subsubsection}{\normalsize\bfseries\color{Navy}}{\thesubsubsection}{.6em}{}
\setlist[enumerate]{leftmargin=2em,itemsep=3pt,topsep=4pt}
\setlist[itemize]{leftmargin=1.6em,itemsep=2pt,topsep=4pt}
\pagestyle{fancy}
\fancyhf{}
\fancyhead[L]{\sffamily\small\color{Muted}COARSE-INFORMATION PROFILES}
\fancyhead[R]{\sffamily\small\color{Muted}RESEARCH DRAFT}
\fancyfoot[L]{\small\color{Muted}28 September 2026}
\fancyfoot[R]{\small\color{Muted}\thepage}
\renewcommand{\headrulewidth}{.3pt}
\fancypagestyle{plain}{\fancyhf{}\fancyfoot[L]{\small\color{Muted}28 September 2026}
\fancyfoot[R]{\small\color{Muted}\thepage}\renewcommand{\headrulewidth}{0pt}}
\theoremstyle{plain}
\newtheorem{theorem}{Theorem}[section]
\newtheorem{lemma}[theorem]{Lemma}
\newtheorem{proposition}[theorem]{Proposition}
\newtheorem{corollary}[theorem]{Corollary}
\theoremstyle{definition}
\newtheorem{definition}[theorem]{Definition}
\newtheorem{example}[theorem]{Example}
\newtheorem{question}[theorem]{Question}
\theoremstyle{remark}
\newtheorem{remark}[theorem]{Remark}
\numberwithin{equation}{section}
\newcommand{\N}{\mathbb N}
\newcommand{\Cantor}{2^{\N}}
\newcommand{\leT}{\leq_{\mathrm T}}
\newcommand{\ltT}{<_{\mathrm T}}
\newcommand{\eqT}{\equiv_{\mathrm T}}
\newcommand{\nleT}{\nleq_{\mathrm T}}
\newcommand{\Core}{\operatorname{Core}}
\newcommand{\CD}{\operatorname{CD}}
\newcommand{\Spec}{\operatorname{Spec}}
\newcommand{\Low}{\operatorname{Low}}
\newcommand{\Comp}{\mathrm{Comp}}
\newcommand{\degT}{\deg_{\mathrm T}}
\newcommand{\bitsum}{\mathbin{\triangle}}
\newcommand{\join}{\mathbin{\oplus}}
\newcommand{\rest}{\mathbin{\upharpoonright}}
\newcommand{\nc}{\mathrm{nc}}
\newcommand{\uc}{\mathrm{uc}}
\newcommand{\eps}{\varepsilon}
\newcommand{\pref}{\preceq}
\newcommand{\concat}{\mathbin{{}^{\frown}}}
\newcommand{\I}{\mathcal I}
\newcommand{\Acal}{\mathcal A}
\newcommand{\zero}{\mathbf 0}
\DeclareMathOperator{\rk}{rank}
\setcounter{tocdepth}{1}

\begin{document}
\thispagestyle{plain}
\begin{center}
{\sffamily\small\bfseries\color{Steel}COMPUTABILITY THEORY\quad /\quad PROVEIT RESEARCH EXTENSION}\par
\vspace{10pt}
{\fontsize{25}{30}\selectfont\bfseries\color{Navy}Finite Realization of\\Coarse-Information Profiles}\par
\vspace{8pt}
{\large Countable Turing ideals, generic secret sharing,\\and simultaneously unattained representative spectra}\par
\vspace{10pt}
{\color{Muted}Prepared for Vladimir Reshetnikov with ChatGPT}\par
{\small\color{Muted}28 September 2026}
\end{center}

\begin{abstract}
For a binary sequence $X$, its coarse-information core consists of the sets
computed by every density-zero perturbation of $X$. We prove a finite
realization theorem: an assignment of countable Turing ideals to coalitions
of finitely many participants is realizable as the cores of their oracle
joins exactly when it is monotone and assigns the computable ideal to the
empty coalition. All nonempty coalition classes can simultaneously be made
to have no least Turing-degree representative. Relative to a presentation
$P$ of the ideals, the full join $H$ can satisfy
$P\ltT H\ltT P'$ and $H'\eqT P'$.

The argument first verifies the single-ideal construction already sketched
in the repository's target C4, then combines robust column coding with
finite additive sharing and a generic-perturbation lemma. A coalition normal
form proves the essential upper bound: its core contains no unintended
information. As a particularly sharp consequence, two sequences with
computable cores can have any prescribed countable Turing ideal as the core
of their join. A separate access-structure version makes every unauthorized
coalition itself $1$-generic relative to the secret presentation.

The non-elementary inputs are named published theorems of Hirschfeldt,
Jockusch, Kuyper, and Schupp. The finite-profile results are presented with
conventional proofs as proposed research contributions, not as externally
certified priority claims or machine-verified theorems. A proof audit,
finite verification program, and further research agenda accompany the text.
\end{abstract}

\noindent\textbf{Scope.} This article does not claim a new solution to the
already-retired least-representative question C1, nor a solution to the
repository's question about \emph{minimal} elements of coarse spectra.
Its main object is simultaneous control of all finite coalition cores.

\newpage
\tableofcontents
\vspace{12pt}
\section*{Reading guide}
For the classification statement, read Section~\ref{sec:main}. For the
shortest substantial example, read the two-share construction in
Section~\ref{sec:two}. Its complete upper-bound mechanism is the transfer
lemma in Section~\ref{sec:generic}.

The proof of the general theorem is concentrated in
Sections~\ref{sec:algebra}--\ref{sec:bounds}: finite sharing, the coalition
normal form, elimination of unintended information, nonattainment, and the
jump calculation. Section~\ref{sec:access} gives the stronger genericity
property for unauthorized coalitions. Section~\ref{sec:attainment} explains
why positive leastness flags cannot be prescribed without further
conditions.

The source and dependency distinctions are in
Section~\ref{sec:provenance}; the audit and finite-check results are in
Section~\ref{sec:audit}. The questions in Section~\ref{sec:questions} are
separate research targets, not additional assertions established by the
main theorem.
\newpage

\section{The problem and the main theorem}\label{sec:main}

\subsection{Information surviving arbitrary sparse errors}

Think of a set of natural numbers as an infinite binary oracle. Its exact
Turing degree measures everything obtainable from a completely accurate
copy. A different question is what remains obtainable when an adversary may
change the oracle on an unknown set of asymptotic density zero. The reduction
used to recover a particular set is allowed to depend on the corrupted
copy. This nonuniformity is essential.

For $E\subseteq\N$ and $N\geq1$, put
\[
 \rho_N(E)=\frac{|E\cap[0,N)|}{N},\qquad
 \delta(X,Y)=\limsup_{N\to\infty}\rho_N(X\bitsum Y).
\]
Here $X\bitsum Y$ is symmetric difference, equivalently bitwise exclusive
or. It is not the oracle join $X\join Y$. Write $D\approx X$ when
$\delta(D,X)=0$, and let
\[
 \CD(X)=\{D\in\Cantor:D\approx X\},\qquad
 \Core(X)=\{A:(\forall D\in\CD(X))\ A\leT D\}.
\]
This is the invariant denoted $X^{\mathfrak c}$ in
Hirschfeldt--Jockusch--Kuyper--Schupp \cite{HJKS}. A \emph{Turing ideal} is a
nonempty collection of sets closed downward under Turing reducibility and
under finite oracle joins. Let $\Comp$ be the ideal of computable sets and
$\Low(B)=\{A:A\leT B\}$ the principal ideal generated by $B$.

For participants $[n]=\{1,\ldots,n\}$, set
\[
 X_C=\bigoplus_{i\in C}X_i\quad(C\ne\varnothing),\qquad
 X_\varnothing=\varnothing.
\]
All finite joins in this article use equal-density periodic interleaving,
in a fixed ordering of their coordinates. In particular, $X_C$ is a
specified oracle, not merely an unspecified representative of a join degree.

\begin{theorem}[Finite profile realization]\label{thm:main}
Let $n\geq1$. An assignment $(\I_C)_{C\subseteq[n]}$ is of the form
\begin{equation}\label{eq:profile}
                \I_C=\Core(X_C)\qquad(C\subseteq[n]).
\end{equation}
for some binary sequences $X_1,\ldots,X_n$ if and only if:
\begin{enumerate}[label=\textup{(\roman*)}]
\item every $\I_C$ is a countable Turing ideal;
\item $\I_\varnothing=\Comp$;
\item $\I_C\subseteq\I_D$ whenever $C\subseteq D$.
\end{enumerate}
Moreover, the realizing sequences can be chosen so that every nonempty
$C$ has \emph{no least Turing degree} in its density-zero class, its uniform
coarse-equivalence class, or its nonuniform coarse-equivalence class.

Choose sequences $(A_{S,k})_{k\in\N}$ generating $\I_S$ for all nonempty
$S\subseteq[n]$, and let $P$ be their combined oracle. There is a realization
for which, with $H=X_{[n]}$,
\begin{equation}\label{eq:jumpbound}
      P\ltT H\ltT P',\qquad H'\eqT P'.
\end{equation}
In particular each $X_i\leT P'$.
\end{theorem}

The assumptions place no further restriction on how much information can
appear when participants cooperate. The ideals for overlapping coalitions
need not obey an intersection identity, and the ideal assigned to a union
may be strictly larger than the ideal generated by the two constituent
ideals. The theorem asserts a complete finite classification at the level
of \emph{cores}, not of the full degree spectra or of coarse reducibility.

\begin{corollary}[Arbitrary two-oracle synergy]\label{cor:synergyintro}
For every countable Turing ideal $\I$, there are $X,Y$ such that
\[
 \Core(X)=\Core(Y)=\Comp,\qquad \Core(X\join Y)=\I.
\]
They can be chosen individually $1$-generic, with no least representative in
any of the three relevant classes of $X$, $Y$, or $X\join Y$.
\end{corollary}

A direct two-oracle construction, stronger than merely applying the main
theorem, appears in Section~\ref{sec:two}. The main theorem itself is proved
in Sections~\ref{sec:algebra}--\ref{sec:bounds}. The construction is finite
in its sharing architecture but infinite in its oracle coding and genericity.

\subsection{Why the upper bound is the real issue}

To make every authorized coalition recover a chosen ideal, it is enough to
supply redundant encodings of its generators. That establishes only one
inclusion in~\eqref{eq:profile}. It does not prevent masks, combinations of
shares, or the choice of the presentation from encoding additional robust
information. Finite probabilistic secrecy by itself also does not establish
an equality of Turing ideals.

Our proof handles this through three levels. The code for an ideal can be
approximated in density by codes computable from finitely many generators.
For a fixed coalition, every partly observed sharing component can be
absorbed into a change of generic coordinates. Finally, a cone-avoidance
compactness theorem converts information surviving all sparse errors into
information surviving a whole positive-radius neighborhood. This lets us
replace each full secret code by a finite approximation and eliminate the
generic coordinates. The surviving information must belong to the prescribed
ideal.

\section{Repository provenance and dependency boundary}\label{sec:provenance}

\subsection{What is already in ProveIt}

The repository was inspected at commit
\begin{center}\small\ttfamily
1a1396d4d3a2ac6812692df520517d86ba3a4785.
\end{center}
The relevant program is
\path{Computability/TuringDegrees/Research/CoarseDegrees/}.
Its research plan already states target C4 and gives the precise dyadic
column construction for realizing one countable Turing ideal as a core
\cite{RepoPlan}. We verify that sketch in Section~\ref{sec:ideal}; we do not
claim it as a new discovery. The same column-coding architecture occurs in
the proof of Theorem~3.11(2) of Hirschfeldt--Jockusch--Schupp
\cite{HJS}, together with their coarse-description projection lemma.

The repository's synthesis already develops equality of representative
spectra, criteria for a least representative, and examples with prescribed
unattained principal infima \cite{RepoSynthesis}. The no-least statement
below is therefore not advertised in isolation as a new phenomenon. Its
role here is that it can be imposed \emph{simultaneously} on every nonempty
coalition while its entire ideal-valued profile is prescribed.

The finite access-structure architecture is related to the classical
construction of general secret-sharing schemes by Ito, Saito, and Nishizeki
\cite{ISN}. We prove the specific additive algebra needed here directly.
No novelty is claimed for splitting one bit into exclusive-or shares. The
proposed contribution is the combination with exact coarse-core upper
bounds, arbitrary countable ideals, and a common relative jump bound.

\subsection{Claim ledger}

\begin{center}\small
\begin{tabularx}{\textwidth}{@{}p{.27\textwidth}Xp{.25\textwidth}@{}}
\toprule
Statement & Role in this article & Status / attribution\\
\midrule
Cone-avoiding compactness; generic trivial core & Two non-elementary inputs & Published, \cite{HJKS}\\
Single-ideal column code & Complete verification of C4 & Already sketched in \cite{RepoPlan}; coding in \cite{HJS}\\
Finite affine sharing & Algebraic mechanism & Classical architecture; direct proof supplied\\
All finite monotone profiles & Main classification theorem & Proposed contribution; full conventional proof\\
Simultaneous nonattainment and relative lowness & Strengthening of realization & Proved here for the constructed profile\\
Unauthorized coalitions generic over the secret presentation & Access-structure refinement & Proved here\\
Uniform exact core is computable & Clarification of C4's uniformity question & Elementary established obstruction, not a novelty claim\\
\bottomrule
\end{tabularx}
\end{center}

A search and source comparison did not establish priority for the main
combined statement. The inspected sources and the repository's stated
results do not supply that statement in this form. This is a bounded
literature assessment, not evidence that no equivalent theorem exists.
Independent mathematical review and a broader priority search are still
appropriate before publication.

\subsection{Formal verification is a separate claim}

The repository distinguishes its genuine oracle-recursion development from
a coarse-degree library containing explicitly admitted literature inputs
\cite{RepoReadme}. This article does not rely on an uninspected assertion
that the relevant Lean declarations are assumption-free. Its proofs are
ordinary mathematical proofs with the two published inputs isolated below.
No Lean or Rocq build was performed for this article. The accompanying
Python program checks finite identities only; its results are not a
substitute for a proof of an infinite computability theorem.

\section{Basic density and oracle calculus}\label{sec:calculus}

\subsection{Periodic coordinates and finite operations}

For $m\geq1$ and $Y_0,\ldots,Y_{m-1}\in\Cantor$, write
\[
 \operatorname{Join}_m(Y_0,\ldots,Y_{m-1})(mt+j)=Y_j(t)
 \quad(0\leq j<m).
\]
There is a computable coordinate extraction map for every $j$. For sets
$E_j$ on the coordinates, counting at lengths divisible by $m$ gives
\begin{equation}\label{eq:joincount}
 \rho_{mN}\!\left(\operatorname{Join}_m(E_0,\ldots,E_{m-1})\right)
       =\frac1m\sum_{j<m}\rho_N(E_j).
\end{equation}
At other lengths the difference is bounded by a term tending to zero.

\begin{lemma}[Density calculus]\label{lem:density}
The following facts hold.
\begin{enumerate}[label=\textup{(\alph*)}]
\item If $D\approx\operatorname{Join}_m(Y_j)$, then every extracted
coordinate $D_j$ satisfies $D_j\approx Y_j$.
\item If $D_j\approx Y_j$ for all $j<m$, then
$\operatorname{Join}_m(D_j)\approx\operatorname{Join}_m(Y_j)$.
\item For two tuples,
\[
 \delta\!\left(\operatorname{Join}_m(Y_j),\operatorname{Join}_m(Z_j)\right)
       \leq\frac1m\sum_{j<m}\delta(Y_j,Z_j).
\]
\item For finite exclusive-or sums,
\[
 \delta\!\left(\mathop{\triangle}_{j<r}Y_j,
                    \mathop{\triangle}_{j<r}Z_j\right)
       \leq\sum_{j<r}\delta(Y_j,Z_j).
\]
\end{enumerate}
Consequently, finite coordinate projections, permutations, duplications,
and pointwise exclusive-or operations preserve coarse similarity whenever
all their input coordinates are coarsely similar.
\end{lemma}
\begin{proof}
For (a), an error in coordinate $j$ before its $N$th position is an error
in the full oracle before position $mN$. Hence its error fraction is at most
$m\rho_{mN}(D\bitsum\operatorname{Join}_m(Y_j))$, which tends to zero.
Parts (b) and (c) follow from~\eqref{eq:joincount} and finite subadditivity
of the limsup. For (d), if the two exclusive-or outputs differ at a position,
at least one corresponding pair of inputs differs there. Thus the output
error set is contained in the finite union of the input error sets. The
last assertion follows by composition of these operations. A fixed extra
oracle used identically on both sides introduces no errors.
\end{proof}

A permutation of finitely many periodic coordinates preserves $\delta$
exactly: at lengths $mN$ it preserves the error count. Regrouping a finite
join of equally sized tuples into one tuple has the same effect. Later each
participant receives exactly the same number of streams, so flattening
coalition joins has no hidden density distortion.

\begin{lemma}[Necessary profile conditions]\label{lem:necessity}
For every $X$, $\Core(X)$ is a countable Turing ideal. If $C\subseteq D$,
then $\Core(X_C)\subseteq\Core(X_D)$. Also
$\Core(X_\varnothing)=\Comp$.
\end{lemma}
\begin{proof}
Every lower cone is downward closed and closed under finite joins. The core
is an intersection of such cones and contains all computable sets. Since
$X$ itself is a coarse description, $\Core(X)\subseteq\Low(X)$; the latter
is countable because there are countably many oracle programs. Given a
coarse description of $X_D$, extract and regroup the coordinates belonging
to $C$. Lemma~\ref{lem:density} produces a coarse description of $X_C$,
proving monotonicity. The empty coalition is represented by a computable
oracle, which proves the final assertion.
\end{proof}

\subsection{Relative cores and the two published inputs}

For a background oracle $B$, define
\[
 \Core^B(X)=\{A:(\forall D\approx X)\ A\leT B\join D\}.
\]
All descriptions here are arbitrary binary sets; they are not required to
be computable in the background oracle.

\begin{theorem}[Cone-avoiding compactness, published input]\label{thm:compact}
If for every $\eps>0$ there is a $Y$ with $\delta(X,Y)\leq\eps$ and
$A\nleT B\join Y$, then there is a $D\approx X$ with
$A\nleT B\join D$.
\end{theorem}

This is Theorem~3.7 of \cite{HJKS}, with all oracle computations relativized
to $B$. Its contrapositive supplies the form we use.

\begin{corollary}[Robust radius]\label{cor:radius}
For all $A,B,X$,
\begin{equation}\label{eq:radius}
 A\in\Core^B(X)
 \quad\Longleftrightarrow\quad
 (\exists\eta>0)(\forall Y)
       [\delta(X,Y)<\eta\ \Longrightarrow\ A\leT B\join Y].
\end{equation}
\end{corollary}
\begin{proof}
If no such radius exists, for every positive $\eps$ choose a counterexample
at distance less than $\eps$ and apply Theorem~\ref{thm:compact}.
Conversely, every coarse description has distance zero and hence is within
any positive radius. The strict inequality avoids irrelevant boundary
conventions; replacing a radius by a smaller one would give the same use.
\end{proof}

\begin{definition}
A sequence $G$ is \emph{$1$-generic relative to $B$} if for every
$B$-computably enumerable set $W$ of finite binary strings, either some
prefix of $G$ belongs to $W$, or some prefix of $G$ has no extension in $W$.
\end{definition}

\begin{theorem}[Generic core, published input]\label{thm:genericcore}
If $G$ is $1$-generic relative to $B$, then
\begin{equation}\label{eq:genericcore}
                     \Core^B(G)=\Low(B).
\end{equation}
\end{theorem}

This is the relativization of Theorem~4.2 of \cite{HJKS}; the background
oracle is present in every recovery computation. We use this precise
relative formulation, not an inference that ordinary genericity is enough
in the presence of an arbitrary background oracle. These two named inputs
are the only non-elementary coarse-computability theorems used below.

\section{Generic coordinates and a transfer principle}\label{sec:generic}

\subsection{The elementary genericity facts needed later}

We record the details to identify precisely which genericity properties the
construction uses. Finite tuples of binary streams are always identified
with one binary sequence by periodic interleaving.

\begin{lemma}[Changes and projections of generic coordinates]\label{lem:genericmaps}
Let $G$ be $1$-generic relative to $P$.
\begin{enumerate}[label=\textup{(\alph*)}]
\item If $B\leT P$, then $G$ is $1$-generic relative to $B$.
\item If $h$ is a $P$-computable homeomorphism of a finite product of
Cantor spaces with $P$-computable inverse, then $h(G)$ is $1$-generic
relative to $P$.
\item Every nonempty finite coordinate projection of $G$ is $1$-generic
relative to $P$.
\end{enumerate}
In particular, invertible affine changes over $\mathbb F_2$, applied
pointwise to finitely many streams with $P$-computable offsets, preserve
$1$-genericity relative to $P$.
\end{lemma}
\begin{proof}
Part (a) follows because every $B$-c.e. set of strings is $P$-c.e. For (b),
view a c.e. set of strings as the effectively open union of its cylinders.
The inverse image under $h$ is $P$-effectively open. If $G$ belongs to it,
then $h(G)$ meets the original open set. Otherwise genericity gives a
cylinder about $G$ disjoint from the inverse image. Its image is an open
neighborhood of $h(G)$ disjoint from the original open set; it contains a
cylinder about $h(G)$, as required.

For (c), pull back the effectively open set by the coordinate projection.
If a cylinder about $G$ avoids this pullback, its projection is an open
neighborhood of the projected sequence and avoids the original set. The
projection of a basic finite-coordinate cylinder is again open. This proves
the avoiding alternative. The meeting alternative is immediate. Finally,
the stated affine maps and their inverses are pointwise finite Boolean
operations, so they satisfy (b).
\end{proof}

\begin{lemma}[Genericity excludes a background-computable description]\label{lem:nodescription}
If $G$ is $1$-generic relative to $P$ and $E\leT P$, then
\[
                         \delta(G,E)=1.
\]
In particular, $G$ has no $P$-computable coarse description.
\end{lemma}
\begin{proof}
Fix a rational $q<1$ and an integer $b$. Consider strings $\sigma$ having a
length $N>b$ with $\rho_N(\sigma\bitsum E)>q$. They form a $P$-c.e.
upward-closed family, where membership is witnessed by a finite prefix.
It is dense: extend any string by a sufficiently long interval on which
the bits are the complements of the corresponding bits of $E$.
A generic sequence must meet a dense c.e. family, since the avoiding
alternative is impossible. Thus there are arbitrarily large such $N$ for
every rational $q<1$. The limsup is one.
\end{proof}

\begin{lemma}[No noncomputable base]\label{lem:nobase}
If $G$ is $1$-generic relative to $P$, $A\leT P$, and $A\leT G$, then
$A$ is computable.
\end{lemma}
\begin{proof}
Choose $\Phi$ with $\Phi^G=A$. The set of strings $\sigma$ for which
$\Phi^\sigma(t)\downarrow\ne A(t)$ for some $t$ is $P$-c.e.
No prefix of $G$ belongs to it. Hence a prefix $\tau$ of $G$ has no
extension in this set. To compute $A(t)$ without an oracle, search finite
extensions of this fixed $\tau$ for any convergent computation
$\Phi^\sigma(t)$. Such an extension exists because $\Phi^G$ is total.
Its output must equal $A(t)$ by the choice of $\tau$. Hard-coding the
finite string $\tau$ gives an ordinary program for $A$.
\end{proof}

\begin{lemma}[A bounded generic with a controlled jump]\label{lem:lowgeneric}
For every $P$ there is a $G\leT P'$ that is $1$-generic relative to $P$.
Every such $G$ satisfies
\begin{equation}\label{eq:genericjump}
 (P\join G)'\eqT G\join P'\eqT P'.
\end{equation}
Consequently $P\ltT P\join G\ltT P'$.
\end{lemma}
\begin{proof}
List the $P$-c.e. families $W_e$ of finite strings. Starting with an empty
string $\sigma_e$, use $P'$ to decide whether there is an extension in
$W_e$. If one exists, find one by $P$-effective search and extend to it;
if none exists, leave the current string unchanged at this substep.
Append a bit to ensure unbounded length. The union $G$ is $P'$-computable
and meets or avoids each $W_e$.

For the jump assertion, fix an index $e$ for a computation with oracle
$P\join G$. The family of finite $G$-strings on which this computation
halts is $P$-c.e. To decide its outcome using $G\join P'$, search prefixes
$\sigma$ of $G$ until either the computation already halts on $\sigma$, or
$P'$ reports that no extension of $\sigma$ permits it to halt.
Genericity guarantees that one alternative eventually occurs. Thus
$(P\join G)'\leT G\join P'$. The reverse reduction uses
$G\leT(P\join G)'$ and monotonicity of the jump. Since $G\leT P'$, the
last degree is $P'$.

Lemma~\ref{lem:nodescription} implies $G\nleT P$, giving the first strict
inequality. If $P\join G\eqT P'$, jumping and using~\eqref{eq:genericjump}
would give $P''\eqT P'$, contrary to strictness of the Turing jump.
\end{proof}

The bound is relative to the \emph{chosen presentation oracle}. A countable
ideal may have many presentations of different degrees. No canonical least
presentation is assumed, and no absolute lowness conclusion is asserted
unless $P$ itself is computable.

\subsection{Generic perturbations of successively simpler data}

The following lemma packages the upper-bound argument. Its strength comes
from quantifying over \emph{every} coarse description of a generic oracle,
not merely the chosen exact generic sequence.

\begin{lemma}[Generic-perturbation transfer]\label{lem:transfer}
Suppose $W$ is a binary sequence and, for each $m$, there are:
\begin{enumerate}[label=\textup{(\roman*)}]
\item an oracle $B_m$ and a sequence $G_m$ that is $1$-generic relative to
$B_m$;
\item a total $B_m$-computable oracle map $F_m:\Cantor\to\Cantor$ such that
$D\approx G_m$ implies $F_m(D)\approx F_m(G_m)$;
\item a bound $\eps_m\to0$ with
$\delta(W,F_m(G_m))\leq\eps_m$.
\end{enumerate}
Then each $A\in\Core(W)$ satisfies $A\leT B_m$ for every sufficiently
large $m$. In particular, if the lower cones $\Low(B_m)$ are increasing,
\begin{equation}\label{eq:transfer}
                 \Core(W)\subseteq\bigcup_m\Low(B_m).
\end{equation}
\end{lemma}
\begin{proof}
Let $A\in\Core(W)$ and take a radius $\eta>0$ from
Corollary~\ref{cor:radius}, with computable background. For every large $m$
we have $\eps_m<\eta$. Fix such an $m$ and any $D\approx G_m$.
The density triangle inequality gives
\[
 \delta(W,F_m(D))
 \leq\delta(W,F_m(G_m))+\delta(F_m(G_m),F_m(D))
 \leq\eps_m<\eta.
\]
Therefore $A\leT F_m(D)\leT B_m\join D$. As $D$ was arbitrary,
$A\in\Core^{B_m}(G_m)$. Theorem~\ref{thm:genericcore} gives $A\leT B_m$.
The same argument works for every sufficiently large $m$. If the lower
cones are increasing, their eventual-membership union is simply the union
in~\eqref{eq:transfer}.
\end{proof}

There is no assumption that $A$ is computable from the master presentation
used to construct the objects. That fact is obtained as a conclusion when
the $B_m$ belong to the intended ideal. This prevents a circular
cone-avoidance argument in which the object to be excluded is silently
placed into the background oracle.

\section{Verifying C4: one countable ideal as a core}\label{sec:ideal}

\subsection{Robust blocks and positive-density columns}

For $A\in\Cantor$, define
\begin{equation}\label{eq:block}
 J(A)(0)=0,\qquad J(A)(t)=A(j)\quad\text{if }2^j\leq t<2^{j+1}.
\end{equation}
The map is uniform in $A$, and $J(A)\eqT A$ by reading the positions $2^j$.

\begin{lemma}[Nonuniform block decoding]\label{lem:block}
Every coarse description $D$ of $J(A)$ computes $A$.
\end{lemma}
\begin{proof}
On the block $[2^j,2^{j+1})$, compute the majority value of $D$, assigning
zero in a tie. If this value is wrong, at least $2^{j-1}$ positions of the
block are errors. Thus the error density at the endpoint $2^{j+1}$ is at
least $1/4$. This can happen for only finitely many $j$ when $D\approx J(A)$.
The majority sequence, computable from $D$, is therefore a finite variant
of $A$. Hard-code its finitely many exceptions to obtain a reduction of
$A$ to $D$.
\end{proof}

The finite correction depends on $D$. We have proved
$\forall D\,\exists e$, not $\exists e\,\forall D$. Neither a bound on the
number of exceptional blocks nor a procedure locating them is part of the
conclusion.

Partition $\N$ into the computable columns
\begin{equation}\label{eq:columns}
 R_k=\{r_k(t):t\in\N\},\qquad
 r_k(t)=2^k(2t+1)-1.
\end{equation}
The partition is disjoint and exhaustive because $k$ is the exponent of two
in $x+1$. The column $R_k$ has density $2^{-(k+1)}$. More precisely,
\begin{equation}\label{eq:tailcount}
 \left|[0,N)\cap\bigcup_{k\geq m}R_k\right|
                  =\left\lfloor\frac{N}{2^m}\right\rfloor.
\end{equation}
Indeed, membership in this union means that $x+1$ is divisible by $2^m$.

Given a sequence $(A_k)$, define its ideal code $Z$ by
\begin{equation}\label{eq:idealcode}
                         Z(r_k(t))=J(A_k)(t).
\end{equation}
Let $Z^{[m]}$ agree with $Z$ on $R_0,\ldots,R_{m-1}$ and be zero on every
other column. Then
\begin{equation}\label{eq:truncation}
 Z^{[m]}\leT B_m:=\bigoplus_{k<m}A_k,
 \qquad \delta(Z,Z^{[m]})\leq2^{-m}.
\end{equation}
The empty join $B_0$ is computable.

\begin{theorem}[C4, with the sketch fully verified]\label{thm:ideal}
The code in~\eqref{eq:idealcode} satisfies
\begin{equation}\label{eq:idealidentity}
                    \Core(Z)=\bigcup_m\Low(B_m).
\end{equation}
Consequently every countable Turing ideal is the core of a binary sequence.
Moreover, the exact code $Z$ is Turing equivalent, uniformly, to the full
presentation $P_A=\bigoplus_{k\in\N}A_k$.
\end{theorem}
\begin{proof}
Let $D\approx Z$ and write $E=D\bitsum Z$. Extract the $k$th column as
$D_k(t)=D(r_k(t))$. Its first $N$ positions lie below $2^{k+1}N$, so
\[
 \rho_N(D_k\bitsum J(A_k))
       \leq2^{k+1}\rho_{2^{k+1}N}(E)\longrightarrow0.
\]
Thus $D_k$ is a coarse description of $J(A_k)$ and computes $A_k$ by
Lemma~\ref{lem:block}. A finite join of these reductions shows
$B_m\leT D$ for every $m$. This proves the inclusion from right to left
in~\eqref{eq:idealidentity}.

For the reverse inclusion, let $A\in\Core(Z)$ and choose a robust radius
$\eta$ from Corollary~\ref{cor:radius}. Choose $m$ with $2^{-m}<\eta$.
Then~\eqref{eq:truncation} gives
$A\leT Z^{[m]}\leT B_m$. This proves the equality. If $\I$ is a countable
ideal, enumerate a generating sequence from it. Its closure under finite
joins and downward reducibility is exactly the right side of the equality.

Finally, $P_A$ computes $Z$ directly from~\eqref{eq:idealcode}. Conversely,
to obtain $A_k(j)$ from the exact code, read
$Z(r_k(2^j))$. This recovers the two-parameter array uniformly, proving the
last assertion.
\end{proof}

\subsection{Two quantifier distinctions}

Every coarse description of $Z$ computes each generator $A_k$ separately.
It need not compute their infinite join $P_A$. The reductions to $A_k$ may
require unrelated finite corrections, with no computable procedure choosing
the corrections as a function of $k$. The exact equivalence $Z\eqT P_A$
is therefore consistent with a nonprincipal core strictly below $\Low(P_A)$.

The theorem also does not say that every degree which is an upper bound for
all the $A_k$ computes a coarse description of $Z$. An arbitrary family of
nonuniform upper-bound reductions need not yield a usable approximation to
the full column array. Thus the theorem is an exact \emph{core}
calculation, not an asserted complete spectrum calculation.

\subsection{Uniform exact recovery collapses}

One question attached to C4 concerns replacing its nonuniform quantifier by
a single recovering functional. In that literal exact-output sense the
answer is completely different, for a reason unrelated to the particular
column code.

\begin{proposition}[Uniform exact core]\label{prop:uniformcore}
For every $X,B$, the collection of sets $A$ for which some fixed functional
$\Phi$ satisfies
\[
                  (\forall D\approx X)\quad\Phi^{B\join D}=A
\]
is exactly $\Low(B)$. In particular the unrelativized uniform exact core of
every $X$ is $\Comp$.
\end{proposition}
\begin{proof}
Suppose such a $\Phi$ exists. For a given input $t$, search with oracle $B$
through all finite binary strings $\sigma$ and all finite computations until
$\Phi^{B\join\sigma}(t)$ converges. This search terminates since
$\Phi^{B\join X}(t)$ does. Every finite $\sigma$ has an extension $D$ that
is a finite variant of $X$: replace the corresponding prefix of $X$ by
$\sigma$. Hence $D\approx X$, and the convergent finite computation must
have output $A(t)$. The search computes $A$ in $B$. Conversely, when
$A\leT B$, simply ignore the $D$-coordinate.
\end{proof}

This is an exact-recovery obstruction, also reflected in the discussion of
uniform decoding in \cite{HJKS}. It is not a statement that all
\emph{uniform coarse reductions} are trivial: such a reduction only has to
output some coarse description of its target, and different input
oracles may produce different output sequences.

\section{Two masks reveal an arbitrary ideal}\label{sec:two}

The smallest instance already captures the reason why cores do not determine
how information combines.

\begin{theorem}[Two-share realization]\label{thm:two}
Let $\I$ be a countable Turing ideal and let $P$ present a generating
sequence. There are binary $X,Y$ such that
\begin{equation}\label{eq:two}
 \Core(X)=\Core(Y)=\Comp,\qquad \Core(X\join Y)=\I.
\end{equation}
Both $X$ and $Y$ are $1$-generic relative to $P$. They and their join have
no $P$-computable coarse descriptions. With a suitable choice, their full
join $H=X\join Y$ satisfies~\eqref{eq:jumpbound}.
\end{theorem}
\begin{proof}
Let $Z$ be the code of Theorem~\ref{thm:ideal}, so $Z\eqT P$. Choose $R$
$1$-generic relative to $P$, and put
\[
                         X=R,\qquad Y=R\bitsum Z.
\]
Translation by the $P$-computable sequence $Z$ is an invertible
$P$-computable homeomorphism. Thus both sequences are $1$-generic relative
to $P$, hence also without an oracle. Theorem~\ref{thm:genericcore} with
computable background gives their two trivial cores.

If $D\approx X\join Y$, its two coordinates $D_0,D_1$ satisfy
$D_0\approx R$ and $D_1\approx R\bitsum Z$. Therefore
$D_0\bitsum D_1\approx Z$. Every member of $\I$ is computed by this
coarse description of $Z$, giving $\I\subseteq\Core(X\join Y)$.

For the opposite inclusion, set
\[
 F_m(V)=V\join(V\bitsum Z^{[m]}),\qquad
 B_m=\bigoplus_{k<m}A_k.
\]
This map is $B_m$-computable and preserves coarse similarity by
Lemma~\ref{lem:density}. Since only the second coordinate changes,
\[
 \delta(X\join Y,F_m(R))\leq 2^{-m-1}.
\]
Also $R$ is $1$-generic relative to every $B_m\leT P$.
Lemma~\ref{lem:transfer} now gives
$\Core(X\join Y)\subseteq\bigcup_m\Low(B_m)=\I$.

Neither $X$ nor $Y$ has a $P$-computable coarse description by
Lemma~\ref{lem:nodescription}. Such a description of their join would
project to one of $X$, so it too is impossible. Finally,
$H\eqT R\join Z\eqT R\join P$. Choosing $R\leT P'$ as in
Lemma~\ref{lem:lowgeneric} gives the jump bound and the strict inequalities.
\end{proof}

The assertion about no least representatives will follow from the general
spectrum criterion in Section~\ref{sec:bounds}. Every member of the three
cores just calculated is computable in $P$, whereas none of these targets
has a $P$-computable description.

\begin{corollary}[There is no core-only operation for oracle join]\label{cor:nooperation}
There is no operation $\mathfrak F$ on pairs of countable Turing ideals
satisfying
\[
             \Core(X\join Y)=\mathfrak F(\Core(X),\Core(Y))
             \qquad\text{for all }X,Y.
\]
\end{corollary}
\begin{proof}
The pair of inputs $(\Comp,\Comp)$ can have every countable Turing ideal
as its output by Theorem~\ref{thm:two}, including two different ideals.
\end{proof}

\begin{example}[A principal secret and a nonprincipal secret]
For a noncomputable $A$, take a presentation of $\I=\Low(A)$. The individual
shares have trivial cores, while their join robustly computes exactly the
sets below $A$. Its core has greatest degree $\degT(A)$, but that degree
will not be the degree of any coarse representative.

For a nonprincipal example, let $\I$ be the ideal generated by the finite
Turing jumps $\varnothing^{(k)}$, $k\in\N$. The strictness of the jump
makes this ideal nonprincipal. Each finite jump survives in every coarse
description of the joined shares, but their infinite presentation need not
survive. In this case the representative spectrum has no greatest lower
bound at all; see Corollary~\ref{cor:infima}.
\end{example}

\section{Finite additive sharing and coalition normal forms}\label{sec:algebra}

\subsection{An all-of-$S$ component}

Let $S=\{s_1<\cdots<s_r\}$ be a nonempty set of participants, and let $Z_S$
be a binary secret stream. For $r\geq2$, introduce $r-1$ mask streams
$R_{S,1},\ldots,R_{S,r-1}$. Give participant $s_j$ the stream
\begin{equation}\label{eq:shares}
 W_{S,s_j}=
 \begin{cases}
 R_{S,j},&j<r,\\
 Z_S\bitsum R_{S,1}\bitsum\cdots\bitsum R_{S,r-1},&j=r.
 \end{cases}
\end{equation}
For $r=1$, give its unique participant $W_{S,s_1}=Z_S$.
Always,
\begin{equation}\label{eq:parity}
                        \mathop{\triangle}_{i\in S}W_{S,i}=Z_S.
\end{equation}

This is the usual additive all-of-$r$ sharing identity. The point of the
next lemma is not that the shares have a probability distribution: no
randomness is assumed. The lemma identifies the observations as a
projection of a computable change of generic coordinates.

\begin{lemma}[Proper observations are free coordinates]\label{lem:affine}
Suppose $Z_S\leT P$. Let $O$ be a proper subset of the $r$ share positions
in~\eqref{eq:shares}. The tuple of shares observed in $O$ is a coordinate
projection of a $P$-computable invertible affine change of the $r-1$ masks.
If $O$ is empty this is the empty projection. If $O$ is nonempty and the
mask tuple is $1$-generic relative to $P$, the observed tuple is
$1$-generic relative to $P$.
\end{lemma}
\begin{proof}
If the last share is not observed, the observations are already distinct
raw masks. Otherwise the last share is observed but at least one earlier
mask position, say $j$, is missing. Replace the $j$th mask coordinate by
\begin{equation}\label{eq:maskchange}
 \widetilde R_{S,j}=Z_S\bitsum\mathop{\triangle}_{\ell<r}R_{S,\ell},
 \qquad \widetilde R_{S,\ell}=R_{S,\ell}\quad(\ell\ne j).
\end{equation}
This is invertible: the inverse recovers the omitted mask by
\[
 R_{S,j}=Z_S\bitsum\widetilde R_{S,j}
                \bitsum\mathop{\triangle}_{\ell<r,\,\ell\ne j}
                         \widetilde R_{S,\ell}.
\]
The new $j$th coordinate is precisely the observed last share; all observed
earlier shares are unchanged. Distinct observations therefore become
distinct new mask coordinates. Since $Z_S\leT P$, the map and its inverse
are $P$-computable. Apply Lemma~\ref{lem:genericmaps} for the last assertion.
When $r=1$, a proper observation can only be empty, which has already been
excluded from the genericity assertion.
\end{proof}

\subsection{Assembling the complete finite profile}

Fix data satisfying Theorem~\ref{thm:main}. For each nonempty
$S\subseteq[n]$, choose generators $(A_{S,k})_{k\in\N}$ of $\I_S$ and let
$Z_S$ be their ideal code. Fix a computable pairing convention and put
\begin{equation}\label{eq:presentation}
 P=\bigoplus_{\substack{\varnothing\ne S\subseteq[n]\\ k\in\N}} A_{S,k}.
\end{equation}
In particular $Z_S\leT P$ uniformly in the finite parameter $S$.

Introduce disjoint coordinates for every mask in every all-of-$S$
component. In addition, introduce one \emph{private generic stream} $T_i$
for each participant $i$. Let $G$ be the regular finite join of all these
mask and private coordinates, chosen $1$-generic relative to $P$.
The number of its coordinates is
\begin{equation}\label{eq:mastercount}
 Q=n+\sum_{\varnothing\ne S\subseteq[n]}(|S|-1)
                 =n+(n-2)2^{n-1}+1\geq1.
\end{equation}
The streams are projections of one generic tuple. We do not separately
choose generic masks and then assume that their joint tuple is generic.

Participant $i$ receives its private stream $T_i$ and its share $W_{S,i}$
for every $S$ containing $i$. There are exactly
\begin{equation}\label{eq:width}
                               L=2^{n-1}+1
\end{equation}
streams per participant. Define $X_i$ to be their regular $L$-way join in
a fixed computable order. For a coalition $C$, flatten its join to the
$M_C=|C|L$ observed streams. This differs from any other fixed ordering only
by a finite periodic-coordinate permutation.

The private streams are not needed for the lower inclusion of cores.
Their purpose is to guarantee that every nonempty coalition contains an
unremovable generic obstruction to a $P$-computable coarse description.
They will also ensure that the normal form always has at least one generic
coordinate.

\begin{lemma}[Coalition normal form]\label{lem:normalform}
For every nonempty $C\subseteq[n]$, there is a finite tuple $V_C$ which is
$1$-generic relative to $P$, and a fixed pointwise finite exclusive-or and
interleaving map $F_C$, such that
\begin{equation}\label{eq:normalform}
       X_C=F_C\bigl(V_C,(Z_S)_{\varnothing\ne S\subseteq C}\bigr).
\end{equation}
Here equality uses the fixed flattened coordinate ordering. The map $F_C$
is an ordinary computable map when the displayed secrets are supplied as
additional inputs. No secret $Z_S$ with $S\nsubseteq C$ is an additional
input. The number of free stream coordinates in $V_C$ is
\begin{equation}\label{eq:freecount}
                  q_C=|C|L-(2^{|C|}-1)\geq |C|.
\end{equation}
\end{lemma}
\begin{proof}
Examine the components separately.

If $S\subseteq C$, all its shares are observed. Keep its $|S|-1$ raw masks
as free coordinates. Reconstruct its last share using those masks and
$Z_S$. When $|S|=1$, there is no free coordinate and the observed stream is
just $Z_S$.

If $\varnothing\ne S\cap C\ne S$, apply Lemma~\ref{lem:affine} to the
proper set of observed positions. All observed shares become free
coordinates of an invertible $P$-computable affine change of the component's
masks. Their reconstruction needs no additional secret input: in the new
coordinates the observed streams are copied verbatim. The unauthorized
secret may enter the \emph{change of coordinates}, but after that change it
is not an input to the reconstruction map.

Components with $S\cap C=\varnothing$ make no observations. Private streams
$T_i$ for $i\in C$ are kept as free coordinates. The affine changes for the
different components act on disjoint groups of coordinates, so together
with the identity on private streams they form one $P$-computable
homeomorphism of the master product space. Project its output to all the
chosen free coordinates. Lemma~\ref{lem:genericmaps} proves that the result
$V_C$ is $1$-generic relative to $P$. It is a nonempty projection because
$C$ contains at least one private stream.

There is one parity constraint for each fully observed component, and no
constraint for a proper observation. There are $2^{|C|}-1$ nonempty subsets
of $C$, so the number of free coordinates is the total number $|C|L$ of
observed streams minus this number of constraints. Each constraint belongs
to a different component and removes exactly one coordinate; hence the
constraints are independent. This proves~\eqref{eq:freecount}, including
the lower bound from the private coordinates.
\end{proof}

The distinction between the coordinate change and the reconstruction map
is central. The first is allowed to use the full presentation $P$ because
its task is to establish relative genericity. The second will use only
finitely much \emph{authorized} information. Giving the reconstruction map
access to $P$ would make the core upper bound far too weak.

\subsection{Truncating only authorized secrets}

For $C\ne\varnothing$, define the finite authorized join
\begin{equation}\label{eq:bcm}
 B_{C,m}=\bigoplus_{\substack{\varnothing\ne S\subseteq C\\ k<m}}A_{S,k}.
\end{equation}
Every such $B_{C,m}$ belongs to $\I_C$ by monotonicity and closure under
finite joins. Its lower cone increases with $m$, and
\begin{equation}\label{eq:idealunion}
                   \bigcup_m\Low(B_{C,m})=\I_C.
\end{equation}
For the reverse inclusion in~\eqref{eq:idealunion}, it is enough that the
terms with $S=C$ already generate $\I_C$.

Let $Z_S^{[m]}$ be the first-$m$-column truncation of $Z_S$. Define a map on
the free generic tuple by
\begin{equation}\label{eq:fcm}
 F_{C,m}(V)=F_C\bigl(V,(Z_S^{[m]})_{\varnothing\ne S\subseteq C}\bigr).
\end{equation}
We identify its finite product domain with $\Cantor$. The map is total and
$B_{C,m}$-computable. By Lemma~\ref{lem:density} it preserves coarse
similarity in the variable $V$.

\begin{lemma}[The quantitative approximation bound]\label{lem:approx}
The normal form satisfies
\begin{equation}\label{eq:approx}
 \delta\bigl(X_C,F_{C,m}(V_C)\bigr)
     \leq\frac{2^{|C|}-1}{|C|L}\,2^{-m}\leq2^{-m}.
\end{equation}
\end{lemma}
\begin{proof}
For a fully observed component $S\subseteq C$, substituting
$Z_S^{[m]}$ for $Z_S$ changes only its last reconstructed stream. The
other $|S|-1$ streams are copied masks. The changed last stream differs on
exactly $Z_S\bitsum Z_S^{[m]}$, which has upper density at most $2^{-m}$.
This description includes singleton components, whose sole stream is the
last stream. Properly observed components and private streams do not change.
There are $2^{|C|}-1$ fully observed components among $|C|L$ output streams.
Apply Lemma~\ref{lem:density}(c). The coefficient is at most one because
each fully observed component contributes a distinct output stream.
\end{proof}

There is no appeal here to cancellation among different error sets. Each
changed stream has its own deterministic upper bound, and only finitely
many streams are involved. This is why a single truncation level works for
all authorized components of a fixed finite coalition.

\section{Proof of simultaneous realization}\label{sec:realization}

\begin{proposition}[Every intended generator survives]\label{prop:lower}
For the construction above, $\I_C\subseteq\Core(X_C)$ for every nonempty
$C\subseteq[n]$.
\end{proposition}
\begin{proof}
Fix a coarse description $D$ of $X_C$. Extract from it the streams belonging
to the all-of-$C$ component. Each extracted stream is a coarse description
of its exact share, by Lemma~\ref{lem:density}(a). The exclusive-or sum of
these descriptions is a coarse description of $Z_C$, by
Lemma~\ref{lem:density}(d) and~\eqref{eq:parity}. Theorem~\ref{thm:ideal}
therefore implies that $D$ computes each member of $\I_C$. As this works
for every $D$, the stated inclusion follows.
\end{proof}

\begin{proposition}[No unintended information survives]\label{prop:upper}
For the same construction, $\Core(X_C)\subseteq\I_C$ for every nonempty
$C\subseteq[n]$.
\end{proposition}
\begin{proof}
Use the fixed generic tuple $V_C$ of Lemma~\ref{lem:normalform}, with
backgrounds $B_{C,m}$ and maps $F_{C,m}$ from~\eqref{eq:fcm}. Since
$B_{C,m}\leT P$, the tuple $V_C$ is $1$-generic relative to $B_{C,m}$.
The maps are total, $B_{C,m}$-computable, and preserve coarse similarity.
The approximation errors tend to zero by Lemma~\ref{lem:approx}.
All hypotheses of Lemma~\ref{lem:transfer} are satisfied. Consequently,
\[
 \Core(X_C)\subseteq\bigcup_m\Low(B_{C,m})=\I_C.
\]
For clarity, the individual quantifier argument is this. Given
$A\in\Core(X_C)$, choose its positive radius $\eta$, then choose $m$ with
$2^{-m}<\eta$. For \emph{every} $D\approx V_C$, the output $F_{C,m}(D)$
lies within that radius. Hence $A\leT B_{C,m}\join D$ for every such $D$.
The relative generic-core theorem forces $A\leT B_{C,m}\in\I_C$.
No index uniformly recovering $A$ from those different outputs is needed.
\end{proof}

\begin{proof}[Proof of the realization equivalence in Theorem~\ref{thm:main}]
Necessity is Lemma~\ref{lem:necessity}. For sufficiency, choose the generator
presentations, codes, masks, and private streams as in
Section~\ref{sec:algebra}. Propositions~\ref{prop:lower} and~\ref{prop:upper}
give every nonempty equality in~\eqref{eq:profile}. The empty equality holds
by convention. The same single master generic works for all coalitions:
there are only finitely many coordinate changes to consider, and each
preserves genericity by a theorem rather than by an extra construction
requirement. The remaining nonattainment and effectiveness clauses are
proved in the next section.
\end{proof}

\subsection{A fully displayed three-participant example}

For $n=3$, the prescribed data are
\[
 \I_1,\I_2,\I_3,\I_{12},\I_{13},\I_{23},\I_{123},
\]
subject only to the inclusions forced by subscripts. Let $Z_S$ code $\I_S$.
Use masks $R_{12},R_{13},R_{23},U,V$ and private streams $T_1,T_2,T_3$,
whose eight-way join is $1$-generic over $P$. Each participant has five
streams:
\begin{center}\small
\begin{tabular}{@{}clllll@{}}
\toprule
Participant & Private & Singleton & First pair & Second pair & Triple\\
\midrule
1 & $T_1$ & $Z_1$ & $R_{12}$ & $R_{13}$ & $U$\\
2 & $T_2$ & $Z_2$ & $Z_{12}\bitsum R_{12}$ & $R_{23}$ & $V$\\
3 & $T_3$ & $Z_3$ & $Z_{13}\bitsum R_{13}$ & $Z_{23}\bitsum R_{23}$ & $Z_{123}\bitsum U\bitsum V$\\
\bottomrule
\end{tabular}
\end{center}

Coalition $\{1,2\}$ fully observes the singleton components 1 and 2 and the
pair component 12. Its observations of 13, 23, and 123 are proper, so the
normal form treats them as generic coordinates. There are ten observed
streams, three full-component parity relations, and seven free streams.
Truncating the three authorized codes changes at most three streams,
with density distance at most $(3/10)2^{-m}$. Thus its core is precisely
$\I_{12}$, even when $\I_{123}$ contains substantially more information.

The full coalition observes fifteen streams and seven independent parity
relations, leaving eight free streams, exactly the original master masks
and private streams. Its core is $\I_{123}$. Nothing in this calculation
requires the three pair ideals to be equal or comparable with one another.

For example, one may prescribe all singleton ideals as $\Comp$, prescribe
$\I_{12}=\Low(A)$ and $\I_{13}=\Low(B)$ for arbitrary $A,B$, prescribe
$\I_{23}=\Comp$, and take $\I_{123}$ to be any countable ideal containing
$A\join B$. Every one of these choices is realized simultaneously.

\section{Representative spectra, nonattainment, and jump bounds}\label{sec:bounds}

\subsection{The spectrum criterion, with its short proof}

For binary targets, write $X\leq_{\nc}Y$ if every coarse description of $Y$
computes a coarse description of $X$. Write $X\leq_{\uc}Y$ if one Turing
functional maps every coarse description of $Y$ to a coarse description of
$X$. Mutual reducibility defines the respective equivalences. Let
$\Spec_r(X)$ be the set of ordinary Turing degrees of representatives
coarsely equivalent to $X$, where $r\in\{\nc,\uc\}$. These spectra should
not be confused with the coarse degrees themselves.

\begin{lemma}[Spectrum identity]\label{lem:spectrum}
For every binary $X$,
\begin{align}\label{eq:spectra}
 \Spec_{\nc}(X)=\Spec_{\uc}(X)
 &=\{\degT(D):D\approx X\}\\
 &=\{\mathbf b:(\exists D\approx X)\ \degT(D)\leq\mathbf b\}.
 \nonumber
\end{align}
The common lower cone of any of these spectra, expressed as a collection
of sets, is $\Core(X)$.
\end{lemma}
\begin{proof}
If $D\approx X$, then $D$ and $X$ have exactly the same coarse descriptions,
so the identity functional gives uniform equivalence. Uniform equivalence
implies nonuniform equivalence. If $Y\equiv_{\nc}X$, apply $X\leq_{\nc}Y$
to $Y$ as a description of itself; it computes a $D\approx X$.
These facts give all inclusions except upward closure of the description
spectrum.

For that closure, let $B$ compute a coarse description $D$ of $X$. Replace
$D(2^k)$ by $B(k)$ for every $k$ and leave all other positions unchanged.
The resulting $E$ still describes $X$ because the powers of two have density
zero. Also $E\leT B$ and $B\leT E$ by reading those powers, so
$\degT(E)=\degT(B)$. This proves~\eqref{eq:spectra}. Intersecting lower cones
over the description spectrum is exactly the definition of the core.
\end{proof}

The lemma is included to make the later inference self-contained; the
same spectrum identity is already in the repository synthesis
\cite{RepoSynthesis}.

\begin{proposition}[When a least representative exists]\label{prop:least}
The common spectrum in~\eqref{eq:spectra} has a least degree if and only if
some member of $\Core(X)$ computes a coarse description of $X$.
\end{proposition}
\begin{proof}
If $A\in\Core(X)$ computes $D\approx X$, then $A\leT D$ by the core
property, so $D\eqT A$. Every member of the common spectrum computes some
description of $X$ and hence computes $A$. Thus $\degT(A)$ is attained and
least. Conversely, let a least degree be represented by $D\approx X$ using
Lemma~\ref{lem:spectrum}. Every other coarse description computes $D$ by
leastness. Thus $D\in\Core(X)$ and $D$ itself is a description.
\end{proof}

\begin{remark}[Allowing natural-number-valued representatives]
These conclusions do not change if the coarse classes contain arbitrary
total functions $\N\to\N$. If the target is binary, convert a description
$E$ to a binary one by assigning 1 exactly where $E$ is 1 and assigning 0
elsewhere. This introduces no error at any previously correct position.
The self-application argument still gives a binary description computable
from every representative. Its degree can be represented by a set coding
the graph of the function. The other inclusions and sparse coding are
already witnessed by binary descriptions. Hence the same spectrum identity
and the same leastness criterion hold in the larger universe.
\end{remark}

\subsection{Nonattainment in every nonempty coalition}

\begin{theorem}[Simultaneous absence of a least representative]\label{thm:noleast}
For the finite-profile construction, every nonempty $C\subseteq[n]$ has no
$P$-computable coarse description and no least degree in any of the three
spectra in~\eqref{eq:spectra}.
\end{theorem}
\begin{proof}
Choose $i\in C$. The private stream $T_i$ is a periodic coordinate of
$X_C$ and is $1$-generic relative to $P$. If a $P$-computable oracle
coarsely described $X_C$, extracting this coordinate would give a
$P$-computable coarse description of $T_i$, contradicting
Lemma~\ref{lem:nodescription}.

Every set in $\Core(X_C)=\I_C$ is computable in $P$: it is reducible to a
finite join of generators from the presentation. Hence no core member can
compute a coarse description of $X_C$. Proposition~\ref{prop:least}
excludes a least degree. The natural-number-valued extension in the previous
remark shows that enlarging the set of representatives cannot restore one.
\end{proof}

\begin{corollary}[Exact behavior of the infimum]\label{cor:infima}
In the constructed profile, fix $C\ne\varnothing$.
If $\I_C=\Low(B)$ is principal, the degree of $B$ is the greatest lower
bound of the representative spectrum but is not in that spectrum.
If $\I_C$ is nonprincipal, the spectrum has no greatest lower bound in the
Turing degrees.
\end{corollary}
\begin{proof}
The common lower cone of the spectrum is exactly $\I_C$ by
Lemma~\ref{lem:spectrum}. A greatest lower bound is precisely a greatest
Turing degree in that ideal. An ideal has such a degree exactly when it is
principal. In the principal case that degree cannot be attained, because
attainment would make it least, contrary to Theorem~\ref{thm:noleast}.
\end{proof}

The statement excludes a \emph{least} element, not necessarily a
\emph{minimal} element. A minimal element need not be below the other
elements. Nothing in the proof identifies the full spectrum closely enough
to exclude minimal elements, and the repository's target C5 is not settled
by this argument.

\subsection{A common effective bound}

\begin{theorem}[Relative lowness of the entire realization]\label{thm:effective}
Choose the master generic $G\leT P'$ as in Lemma~\ref{lem:lowgeneric}.
Then the full join $H=X_{[n]}$ satisfies
\[
                    H\eqT P\join G,
\]
and hence $P\ltT H\ltT P'$ and $H'\eqT P'$.
For every nonempty $C$ there is also the relative bound
\begin{equation}\label{eq:coalitionjump}
 P\ltT P\join X_C\ltT P',\qquad (P\join X_C)'\eqT P'.
\end{equation}
\end{theorem}
\begin{proof}
All shares are computable from $P\join G$, so $H\leT P\join G$.
Conversely, the full join contains every private stream and every raw mask:
for an all-of-$S$ component the first $|S|-1$ shares are exactly its masks.
Thus $H$ computes $G$. Taking the exclusive-or of the complete shares in
each component recovers its exact code $Z_S$. Exact decoding of the columns
of all these codes recovers the full array~\eqref{eq:presentation}
uniformly. Hence $P\leT H$, proving $H\eqT P\join G$.
Lemma~\ref{lem:lowgeneric} gives the stated full-join conclusions.

For a nonempty coalition, $P\join X_C\leT P\join G\eqT H$, so
$(P\join X_C)'\leT H'\eqT P'$. The reverse reduction follows from
$P\leT P\join X_C$. The private stream of any member of $C$ shows
$X_C\nleT P$, proving the lower strict inequality. If the upper one failed,
$P\join X_C\eqT P'$ and jumping would again give $P''\eqT P'$.
\end{proof}

\begin{proof}[Completion of Theorem~\ref{thm:main}]
The realization equivalence was proved in Section~\ref{sec:realization}.
Theorem~\ref{thm:noleast} supplies its simultaneous nonattainment clause,
and Theorem~\ref{thm:effective} supplies~\eqref{eq:jumpbound}.
All clauses refer to the same construction.
\end{proof}

\begin{remark}[What the effective bound does and does not provide]
The argument is uniform at the level of the given presentation: a fixed
$P'$-effective generic construction and fixed finite formulas produce the
shares. This is not a computable algorithm that takes an arbitrary abstract
ideal as a finite input. Also, $X_C$ need not compute $P$, so
$X_C'\eqT P'$ is not claimed for every proper coalition. The precise
uniform statement is~\eqref{eq:coalitionjump}, with the background joined in.
\end{remark}

\section{Access structures with genuinely generic unauthorized views}\label{sec:access}

The general profile theorem permits different nontrivial ideals for many
coalitions. An access structure is the useful special case where the only
desired cores are $\Comp$ and one fixed secret ideal. There is then a
stronger realization: unauthorized views themselves, before any corruption,
can be $1$-generic relative to the entire secret presentation.

\begin{definition}
An \emph{access structure} on $[n]$ is an upward-closed family
$\Acal\subseteq\mathcal P([n])$ with $\varnothing\notin\Acal$. Its members
are called authorized. The empty family is allowed. Let
$\min\Acal$ be its inclusion-minimal members.
\end{definition}

\begin{theorem}[Generic access-structure realization]\label{thm:access}
Let $\I$ be a countable Turing ideal presented by $P$, and let $\Acal$ be a
finite access structure. There are binary participant oracles $X_1,\ldots,X_n$
such that, for every coalition $C$,
\begin{equation}\label{eq:access}
 \Core(X_C)=
 \begin{cases}
 \I,&C\in\Acal,\\
 \Comp,&C\notin\Acal.
 \end{cases}
\end{equation}
For each nonempty unauthorized $C$, the exact oracle $X_C$ is $1$-generic
relative to $P$. In particular,
\begin{equation}\label{eq:exactsecrecy}
                  \Low(X_C)\cap\Low(P)=\Comp.
\end{equation}
Every nonempty coalition class has no least representative. All shares can
be chosen computable in $P'$, with
$(P\join X_C)'\eqT P'$ and $P\ltT P\join X_C\ltT P'$ for nonempty $C$.
When $\Acal$ is nonempty, the full join itself satisfies
$H\eqT P\join G$ and the sharper full-join bound~\eqref{eq:jumpbound}.
\end{theorem}
\begin{proof}
Let $Z$ be the single code of the generators of $\I$, so $Z\eqT P$.
For every $S\in\min\Acal$, create an all-of-$S$ component with the same
secret stream $Z$, using disjoint mask coordinates for different components.
Every participant receives at least one private generic stream.

To keep all joins equally interleaved, let
\[
 d_i=|\{S\in\min\Acal:i\in S\}|,\qquad
 L_*=1+\max_i d_i.
\]
Give participant $i$ its $d_i$ component streams and exactly $L_*-d_i$
additional, distinct private generic streams. These counts are positive.
All mask and private streams form one master tuple $G$, chosen $1$-generic
relative to $P$. Participant $i$'s oracle is their $L_*$-way join.

If $C$ is unauthorized, it contains no $S\in\min\Acal$. Its observations
of every component are therefore proper. By the same affine changes as in
Lemma~\ref{lem:normalform}, its entire flattened view is a nonempty
coordinate projection of a $P$-computable homeomorphic image of $G$.
There are no fully observed secrets to reconstruct. Thus $X_C$ itself is
$1$-generic relative to $P$. Theorem~\ref{thm:genericcore} with computable
background gives its ordinary core $\Comp$, and
Lemma~\ref{lem:nobase} gives~\eqref{eq:exactsecrecy}.

If $C$ is authorized, it contains at least one minimal authorized set $S$.
The corresponding parity sum of coarse share descriptions yields a coarse
description of $Z$. Thus $\I\subseteq\Core(X_C)$. For the reverse
inclusion, form the same coalition normal form, with one copy of $Z$ for
each fully observed component. Replace all those copies by $Z^{[m]}$.
Only finitely many output streams change, each on a set of upper density at
most $2^{-m}$. The reconstruction map is computable in
$B_m=\bigoplus_{k<m}A_k$ and preserves coarse similarity in its generic
coordinates. Lemma~\ref{lem:transfer} gives
$\Core(X_C)\subseteq\bigcup_m\Low(B_m)=\I$.

Every nonempty coalition contains a private generic stream and therefore
has no $P$-computable description. All members of its core belong to
$\Low(P)$, so Proposition~\ref{prop:least} excludes a least degree.
Choose $G\leT P'$ for the bounds. We have $X_C\leT P\join G$, while a
private stream forces $X_C\nleT P$. The proof of
Theorem~\ref{thm:effective} yields the relative jump assertions.
If $\Acal\ne\varnothing$, the full coalition sees a complete component,
so recovers $Z\eqT P$, as well as all masks and private streams; hence
$H\eqT P\join G$. If $\Acal$ is empty, all views are simply joins of
private generic streams, and no recovery of $P$ from $H$ is claimed.
\end{proof}

\subsection{Thresholds and a worked example}

For $1\leq t\leq n$, take $\Acal=\{C:|C|\geq t\}$.
Its minimal members are exactly the $t$-element subsets. Thus any threshold
pattern is realized, and each party receives
$\binom{n-1}{t-1}$ component streams plus one private stream. This count
measures the width of this particular construction, not an optimality claim.

For a two-of-three structure, one can write the whole scheme as
\[
\begin{aligned}
 X_1&=\operatorname{Join}_3(T_1,R_{12},R_{13}),\\
 X_2&=\operatorname{Join}_3(T_2,Z\bitsum R_{12},R_{23}),\\
 X_3&=\operatorname{Join}_3(T_3,Z\bitsum R_{13},Z\bitsum R_{23}).
\end{aligned}
\]
The six-way join of the $T$ and $R$ streams is generic over $P\eqT Z$.
Each singleton view is generic over $P$ and computes no noncomputable
$P$-computable set. Each pair robustly recovers the ideal $\I$ through its
complete pair component, and its core contains nothing else. The triple has
the same core $\I$, even though it observes all three redundant copies.
Every nonempty view lacks a least-degree coarse representative.

\subsection{Why this is not a practical cryptographic claim}

Equation~\eqref{eq:exactsecrecy} is a precise statement about exact Turing
reductions, and~\eqref{eq:access} is a precise statement about information
surviving arbitrary density-zero errors. They do not assert finite-key
security, computational indistinguishability, Shannon secrecy of a sampled
finite distribution, resistance to finite leakage models, or efficient
share generation. The streams may be noncomputable, and the construction
of a suitable generic uses a jump oracle. Its relevance is mathematical:
it separates exact information, robust nonuniform information, and the
finite combinatorics of cooperation.

\section{A boundary: leastness cannot be prescribed arbitrarily}\label{sec:attainment}

The main theorem realizes every monotone ideal profile while making all
nonempty spectra unattained. Replacing that uniform negative requirement
by arbitrary positive and negative leastness flags introduces genuine
compatibility conditions.

\begin{proposition}[Least representatives are closed under finite joining]\label{prop:leastjoin}
Suppose the coarse spectra of $X$ and $Y$ have least degrees $\mathbf a$
and $\mathbf b$. Then the coarse spectrum of $X\join Y$ has least degree
$\mathbf a\join\mathbf b$, and
\[
            \Core(X\join Y)=\Low(A\join B)
\]
for representatives $A,B$ of those two degrees.
\end{proposition}
\begin{proof}
By Proposition~\ref{prop:least}, choose coarse descriptions $D_X\approx X$
and $D_Y\approx Y$ with $D_X\eqT A$ and $D_Y\eqT B$, where
$A\in\Core(X)$ and $B\in\Core(Y)$. Their join is a coarse description of
$X\join Y$ computable from $A\join B$.
Every description of $X\join Y$ projects to descriptions of $X$ and $Y$,
so it computes both $A$ and $B$, hence $A\join B$. The joined description
therefore has least degree $\mathbf a\join\mathbf b$.
The core is at least $\Low(A\join B)$ by this recovery and at most that
ideal because this particular description has exactly that degree.
\end{proof}

For a participant family, let $\mathcal L$ contain exactly the coalitions
whose spectra have a least degree. It contains the empty coalition. If
$C,D\in\mathcal L$, then $C\cup D\in\mathcal L$: joining their views and
then deleting duplicated coordinates is uniformly coarse-equivalent to the
view of the union. Moreover, the least degree at the union must be the join
of the two least degrees. Thus $\mathcal L$ must be union-closed, and its
principal core labels must satisfy this join identity. These conditions are
not imposed on coalitions outside $\mathcal L$.

This explains why an example of arbitrary synergy cannot be obtained by
using individually coarsely computable shares. A finite join of computable
coarse descriptions is again computable. In the two-share construction the
individual cores are computable, but no computable member of either core
computes a description of its share. The gap between having a principal
core and attaining its greatest degree is exactly what makes arbitrary
joint information possible.

\section{Proof audit and reproducible finite verification}\label{sec:audit}

\subsection{Logical dependency chain}

The main proof can be checked in the following order. The only external
coarse-computability inputs are Theorems~\ref{thm:compact}
and~\ref{thm:genericcore}. The first yields the robust-radius corollary.
Finite counting gives the block and column calculations, establishing the
single-ideal identity. The two inputs together yield the transfer lemma.
Finite affine algebra and elementary genericity invariance yield the
coalition normal form. That form and the explicit approximation bound feed
the transfer lemma and give the upper inclusion. Decoding one fully observed
component gives the lower inclusion independently. Private generic streams
then give nonattainment, and exact access to the full array gives the jump
bound.

In particular, the proof of the upper inclusion does not use the main
realization theorem, the desired ideal equality, the absence of a least
representative, or any claim that the output is minimal. There is no
self-reference through the intended conclusion.

\begin{center}\small
\begin{longtable}{@{}p{.29\textwidth}p{.64\textwidth}@{}}
\toprule
Potential failure & Check or repair used here\\
\midrule
\endhead
Coarse information confused with exact information & $Z_S\eqT P_S$ is an exact statement; description decoding is only per-generator and nonuniform.\\
Ordinary genericity used with an unauthorized oracle & The master tuple is generic relative to the entire presentation $P$; the reconstruction step uses only $B_{C,m}$.\\
Separately generic masks assumed jointly generic & All masks and padding streams are projections of one master generic.\\
Unauthorized secrets left in the final reduction oracle & They enter an affine coordinate change, not the truncated reconstruction map.\\
Genericity assumed under arbitrary Turing equivalence & Only explicit homeomorphisms and finite open coordinate projections are used.\\
A coarse description extracted from a zero-density channel & All extraction channels here have fixed positive periodic or dyadic density.\\
An infinite union of zero-density error sets used & Every decoding or perturbation comparison uses only finitely many channels at once. Infinite columns are controlled by the explicit tail density $2^{-m}$.\\
Finite majority correction treated as uniformly available & The correction is hard-coded separately for each description and generator; no uniform exact decoder is asserted.\\
No least element confused with no minimal element & The proof uses the leastness criterion only. Minimal elements remain unclassified.\\
Relative lowness confused with absolute lowness & The exact bounds are $H'\eqT P'$ and $(P\join X_C)'\eqT P'$.\\
Finite testing advertised as theorem verification & The Python tests check only the finite algebra and counting listed below.\\
\bottomrule
\end{longtable}
\end{center}

\subsection{The finite checks actually run}

The archive contains \path{checks/verify_finite.py}, requiring only the
Python standard library. It writes \path{checks/results.json}. The run
included with this article passed all checks in the following table.

\begin{center}\small
\begin{tabularx}{\textwidth}{@{}Xp{.35\textwidth}@{}}
\toprule
Check & Executed scope\\
\midrule
All-of-$r$ exclusive-or reconstruction & $1\leq r\leq8$; 510 cases\\
Proper observations have identical finite distributions for either secret bit & 502 observed-subset comparisons\\
Coalition mask ranks and independent secret constraints & All 247 nonempty coalitions for $1\leq n\leq7$\\
Affine normal-form reconstruction and inverse & 69,888 seeded cases, $n\leq6$\\
Exact dyadic tail formula & 22,539 pairs $(N,m)$\\
Column-index inversion & 50,000 positions\\
Exclusive-or error inclusion & 5,460 exhaustive bit-tuple pairs\\
Majority correctness and necessary error count & 131,628 cases; block length at most 16\\
Recovery of finite monotone access families from minimal members & All 167 four-participant families with unauthorized empty coalition\\
\bottomrule
\end{tabularx}
\end{center}

The rank test represents the masks as independent formal coordinates over
$\mathbb F_2$. For each coalition it checks
\[
 \rk(\text{mask observation matrix})=|C|L-(2^{|C|}-1).
\]
It also checks that adjoining the formal secret coordinates makes the
observation rows independent. Together these tests support the finite
normal-form bookkeeping. They do not test genericity of an infinite
sequence, recovery from every density-zero perturbation, or any oracle
jump equality. Those are proved in the article, using the declared inputs.

The fixed seed makes the non-exhaustive reconstruction checks reproducible;
the distribution, rank, tail, majority, and access-family checks are
exhaustive over their stated finite ranges. No statistical confidence
claim is attached to the seeded checks.

\subsection{A route to a genuine formalization}

A faithful implementation should first define binary coarse descriptions,
periodic joins, and dyadic column extraction over the repository's actual
oracle semantics. The elementary density lemmas and the block decoder can
then be formalized without any genericity machinery. The decoder's output
should be an existential Turing reduction, with its finite correction
visible; a uniform decoder theorem would be false.

The ideal code should be represented by a two-parameter array, keeping the
array oracle separate from the ideal generated by its finite joins. The
finite affine part can be proved independently over vectors over
$\mathbb F_2$, then lifted pointwise to streams. Useful intermediate theorem
interfaces are: invertibility of the missing-mask replacement, preservation
of genericity by an effective homeomorphism and a finite projection, and
coarse-similarity preservation by reconstruction.

The two published inputs can initially be explicit theorem parameters in a
conditional development. Such a development would establish a correct
implication from those inputs but would not be an assumption-free proof of
the final theorem. A complete formalization must either prove the inputs or
import already proved declarations whose dependencies have been audited.
The last layers are the transfer lemma, finite coalition assembly,
nonattainment criterion, and the ordinary generic jump calculation. The
finite code shipped here can generate examples for this effort, but it is
not itself a proof-producing formalizer.

\section{Further research questions and topics}\label{sec:questions}

The questions below are proposed continuations of this article. Unless a
repository label is explicitly mentioned, they are not claimed to be
published open problems. They separate issues actually settled above from
those that require additional ideas.

\subsection{Full proof verification and priority}

\begin{question}[Formalize the finite-profile theorem]\label{q:formal}
Can the complete theorem, including its two published inputs, be added to
ProveIt's genuine oracle-semantics development with an explicit axiom audit
and no admitted mathematical statements?
\end{question}

A useful first milestone is the complete finite affine and density layer;
a second is the single-ideal identity; the critical third is the
relative generic-core and compactness interface. A literature comparison
should run in parallel, searching for equivalent ideal-valued or
mass-problem formulations rather than only the terminology ``coarse core.''
A successful priority search may reclassify the result as a synthesis; it
would not invalidate the construction or its proofs.

\subsection{Countably many participants}

\begin{question}[A coherent countable realization theorem]\label{q:countable}
Given countable Turing ideals assigned monotonically to all finite subsets
of $\N$, with computable empty ideal, can they be realized by one countable
family $(X_i)_{i\in\N}$? Which additional conditions are needed if ideals
are also assigned to infinite coalitions?
\end{question}

Applying the finite theorem independently to initial segments does not
answer this: the resulting shares need not agree on common participants.
A countable construction has infinitely many incident components and no
single finite mask rank. One possible route is to put components into
shrinking positive-density columns and prove a tail approximation lemma
that eliminates all but finitely many observed components. The main danger
is unauthorized information in the infinite collection of coordinate
changes. Infinite coalition joins introduce a further distinction between
uniform access to all participants and individual nonuniform reductions.

\subsection{Extending fixed shares rather than rebuilding them}

\begin{question}[An extension property]\label{q:extension}
Given already fixed $X_1,\ldots,X_n$, which monotone extensions of their
core profile to $n+1$ participants are realizable while preserving those
oracles exactly, or at least preserving their coarse-equivalence classes?
\end{question}

The present theorem rebuilds all shares from one master generic. It does
not assert an amalgamation or extension property for arbitrary existing
shares. Exact preservation may impose spectrum constraints invisible in
the cores. Even the two-share question, with $X$ fixed and only $Y$ free,
asks for more than Theorem~\ref{thm:two}. It would be useful to identify a
class of extendible generic presentations and a precise obstruction outside
that class.

\subsection{Prescribing attainment together with the ideals}

\begin{question}[Mixed leastness profiles]\label{q:attainment}
Characterize the pairs consisting of a finite ideal profile $(\I_C)$ and
a family $\mathcal L$ of coalitions whose spectra must have least degrees.
Are principality on $\mathcal L$, union closure of $\mathcal L$, and the
join identities of Proposition~\ref{prop:leastjoin} sufficient, or are
there further compatibility conditions?
\end{question}

Our construction settles the case $\mathcal L=\{\varnothing\}$. Positive
attainment cannot simply be enforced by adding a least-degree block code:
that may expose information to larger coalitions and can conflict with
prescribed nonattainment there. A promising finite starting case is three
participants with only one singleton and one overlapping pair required to
attain their cores.

\subsection{The full spectra and the repository's C5}

\begin{question}[Beyond common information]\label{q:spectrum}
What are the complete coarse representative spectra of the two-share code
and the finite-profile construction? Can the generic masks be chosen so
that selected spectra have no minimal element, while preserving every core
equality and the common jump bound?
\end{question}

This directly touches C5 of the repository plan \cite{RepoPlan}, but the
present no-least theorem does not answer it. The core gives the common lower
cone, not the collection of degrees which compute a description. Progress
would require an effective description-building or description-lowering
argument, not just a further common-information calculation. A first target
is the two-share case for one principal ideal, where the intended infimum
is completely explicit.

\subsection{Degree bounds independent of a chosen enumeration}

\begin{question}[Presentation-sensitive complexity]\label{q:bounds}
For a fixed abstract countable ideal, what is the least possible complexity
of a presentation admitting the bounds proved here? Can one characterize
when a given oracle $Q$ bounds a realization with the required exact
profile and nonattainment?
\end{question}

The theorem gives a uniform upper bound from $P'$, not an optimal bound
for the ideal itself. Rearranging a generating sequence can add information
to its presentation without changing the ideal. A useful invariant would
separate the cost of presenting the ideal from the cost of producing a
generic extension. Asking to construct these particular nonattained
realizations below the same $P$ is impossible because they intentionally
have no $P$-computable descriptions; meaningful improvements must alter the
presentation bound or the construction requirements.

\subsection{Uniform reductions and other description notions}

\begin{question}[Uniform coarse structure of the profile]\label{q:uniform}
Can one prescribe, along with the core profile, selected uniform coarse
reducibilities among coalition views? Which finite patterns are compatible
with the nonuniform recoveries used here?
\end{question}

Proposition~\ref{prop:uniformcore} settles only literal uniform
\emph{exact} recovery. It says nothing comparable about uniform production
of descriptions. The component parity maps are uniform at the description
level, whereas block correction is not uniform at the exact-set level.
Keeping these two interfaces separate may lead to a useful classification
of uniform coarse access structures. Effective-dense descriptions, which
permit explicit omissions but no wrong answers, require a separate study;
the density-zero error argument cannot be transferred merely by changing
terminology. This is relevant to C2 and C8 of \cite{RepoPlan} without
claiming to resolve them.

\subsection{Efficient finite architectures}

\begin{question}[Stream complexity and linear sharing]\label{q:efficiency}
For a prescribed finite ideal profile, what is the smallest number of
positive-density streams per participant in a realization of this form?
Can the all-subset construction be replaced by a smaller linear sharing
architecture while preserving the generic normal form?
\end{question}

The construction uses $2^{n-1}+1$ streams per participant. That is an
explicit bound, not a lower bound. For access structures, minimal authorized
sets already reduce the number of components. More efficient finite linear
schemes are candidates, but an adaptation must prove that an unauthorized
observation is a projection of an effective homeomorphic change of masks,
and must give a density-continuity estimate for reconstruction. Finite rank
alone should not be substituted for these oracle-level assertions.

\subsection{Other error ideals and higher reducibilities}

\begin{question}[Replacing asymptotic density zero]\label{q:errors}
For which ideals of allowed error sets does a finite monotone-profile
realization theorem hold? What changes for upper Banach density zero,
logarithmic density zero, or a fixed summable ideal?
\end{question}

The proof separates three required ingredients: a robust code on channels
from which descriptions can be extracted; a finite-operation error
calculus and controlled tails; and an analogue of the robust-radius plus
generic-core mechanism. The finite sharing algebra survives unchanged,
but the compactness and coding properties do not follow from abstract
ideal closure. A carefully chosen error ideal may therefore supply either
a generalization or a structural counterexample.

\begin{question}[Hyperarithmetic profiles]\label{q:hyper}
Can arbitrary finite monotone profiles of suitable countable
hyperdegree ideals be realized as information common to all coarse
descriptions, with hyperarithmetic rather than Turing reducibility?
\end{question}

The repository's report 10 discusses a failure of the corresponding
hyperarithmetic robust-radius principle \cite{RepoHyper}. Consequently the
transfer lemma must not simply be reprinted with $\leq_h$ in place of
$\leT$. One needs a replacement for the step which turns recovery from all
coarse descriptions into recovery from nearby truncated reconstructions.
The finite affine normal form may still be useful, but it does not solve
that missing higher-computability problem.

\section{Conclusion}\label{sec:conclusion}

The core records common robust information, not the difficulty of choosing
a representative and not all information present in exact oracles.
The finite realization theorem shows that its interaction with cooperation
is as unrestricted as monotonicity permits: every finite monotone diagram
of countable Turing ideals can be realized, even while every nonempty
representative spectrum fails to attain a least degree.

The key proof is the coalition upper bound. A generic normal form isolates
all fully observed codes; their density truncations require only finitely
many authorized generators. Cone-avoiding compactness and the relative
generic-core theorem then remove the free coordinates. This proves an
exact ideal equality rather than merely showing that some chosen secret
is recoverable. The same construction has a common relative low jump,
and the access-structure variant makes all unauthorized exact views generic
over the entire presentation.

The single-ideal theorem verifies the repository's C4 sketch rather than
superseding its provenance. The proposed new contribution is the
simultaneous finite classification and its refinements. Full spectra,
minimal elements, coherent countable extensions, and mixed attainment
patterns remain separate research targets. The proofs and finite checks
here provide explicit objects and interfaces with which to attack them.

\appendix
\section{Notation and construction index}
\begin{center}\small
\begin{longtable}{@{}p{.26\textwidth}p{.67\textwidth}@{}}
\toprule
Notation & Meaning\\
\midrule
\endhead
$X\bitsum Y$ & Symmetric difference / pointwise exclusive or.\\
$X\join Y$ & Oracle join; throughout, a specified equal-density interleaving.\\
$\delta(X,Y)$ & Upper asymptotic density of the disagreement set.\\
$D\approx X$ & $\delta(D,X)=0$; $D$ is a coarse description of $X$.\\
$\Core^B(X)$ & Sets computed by $B\join D$ for every $D\approx X$.\\
$\Core(X)$ & The unrelativized core, with computable background.\\
$\Low(B)$ & The Turing ideal $\{A:A\leT B\}$.\\
$J(A)$ & Redundant block code, constant on $[2^j,2^{j+1})$.\\
$R_k$ & Positive-density column $\{2^k(2t+1)-1:t\in\N\}$.\\
$Z_S$ & Ideal code of the generator sequence $(A_{S,k})_k$.\\
$Z_S^{[m]}$ & First-$m$-column truncation; the remaining columns are zero.\\
$P$ & Oracle presenting all generator sequences; not assumed to belong to any prescribed ideal.\\
$G$ & Single master tuple of masks and private streams, $1$-generic relative to $P$.\\
$W_{S,i}$ & Participant $i$'s share in the all-of-$S$ component.\\
$T_i$ & A private generic stream ensuring nonattainment.\\
$X_C$ & View of coalition $C$, the finite join of its participant oracles.\\
$V_C$ & Free generic coordinates of the coalition normal form.\\
$B_{C,m}$ & Finite join of generators in authorized components $S\subseteq C$ with $k<m$.\\
$F_{C,m}$ & Reconstruction with authorized secrets truncated to $m$ columns.\\
$H$ & Full join $X_{[n]}$, equivalent to $P\join G$ in the main construction.\\
\bottomrule
\end{longtable}
\end{center}

\section{Source review and reproducibility details}\label{app:sources}

The repository comparison used the pinned commit displayed in
Section~\ref{sec:provenance}, not an unversioned assumption about future
contents. The inspected core documents were the TuringDegrees README,
the coarse research plan's C1--C8 discussion, the synthesis's core and
spectrum sections and result inventory, and report 10's README. The
literature comparison used the HTML text of the two arXiv papers in the
bibliography, the author-hosted publication record for \cite{HJKS}, and
the publisher record for \cite{ISN}. The latter was used for attribution
of the classical access-structure architecture, not as an uninspected
source for any computability theorem.

The proof uses the theorem numbering of the cited arXiv versions. All
newly derived claims are proved in this article. No claim is made that the
repository was exhaustively audited file by file or that the literature
search was exhaustive. Search results with no relevant mathematical
content were not treated as evidence of originality.

The source is a standalone \LaTeX{} file with an embedded bibliography.
The build script runs pdf\LaTeX{} three times to resolve references and
contents. It requires the New PX font packages and standard AMS,
geometry, microtype, booktabs, enumitem, xurl, titlesec, and hyperref
packages available in common TeX Live installations. No font binaries are
included in the archive. Run the finite checks separately with
\texttt{python3 checks/verify\_finite.py}. The supplied PDF was compiled
from this source; document-layout checks are recorded in the accompanying
proof audit rather than confused with mathematical verification.

\begin{thebibliography}{99}

\bibitem{HJKS}
D.~R. Hirschfeldt, C.~G. Jockusch, Jr., R.~Kuyper, and P.~E. Schupp,
\emph{Coarse reducibility and algorithmic randomness},
Journal of Symbolic Logic \textbf{81} (2016), 1028--1046.
Theorem 3.7 (cone-avoiding compactness) and Theorem 4.2 (generic core).
Preprint: \url{https://arxiv.org/abs/1505.01707}.
Text checked: \url{https://arxiv.org/html/1505.01707v1}.
Publication record: \url{https://www.math.uchicago.edu/~drh/Papers/coarsereducibility.html}.

\bibitem{HJS}
D.~R. Hirschfeldt, C.~G. Jockusch, Jr., and P.~E. Schupp,
\emph{Coarse computability, the density metric, Hausdorff distances between
Turing degrees, perfect trees, and reverse mathematics},
arXiv:2106.13118v1 (2021).
Lemma 3.10 and the proof of Theorem 3.11(2).
\url{https://arxiv.org/html/2106.13118v1}.

\bibitem{ISN}
M.~Ito, A.~Saito, and T.~Nishizeki,
\emph{Secret sharing scheme realizing general access structure},
Electronics and Communications in Japan (Part III: Fundamental Electronic
Science) \textbf{72}, no.~9 (1989), 56--64.
\url{https://doi.org/10.1002/ecjc.4430720906}.

\bibitem{RepoPlan}
V.~Reshetnikov, ProveIt repository,
\emph{Turing Degrees: A Unified Research Report}, second edition,
18 September 2026, Section ``Coarse and effective-dense representative
spectra,'' especially target C4 and Proposition ``proposed.''
Repository snapshot
\texttt{1a1396d4d3a2ac6812692df520517d86ba3a4785}.
\url{https://github.com/VladimirReshetnikov/ProveIt/blob/1a1396d4d3a2ac6812692df520517d86ba3a4785/Computability/TuringDegrees/Research/CoarseDegrees/research-plan/turing_degrees_unified.tex}.

\bibitem{RepoSynthesis}
V.~Reshetnikov, ProveIt repository,
\emph{Coarse Classes Without Least Turing Degrees: Deduplicated results of
the nine research reports on target C1}, 18 September 2026.
Same repository snapshot. Conventional research synthesis; status and
provenance distinctions are retained in the present article.
\url{https://github.com/VladimirReshetnikov/ProveIt/blob/1a1396d4d3a2ac6812692df520517d86ba3a4785/Computability/TuringDegrees/Research/CoarseDegrees/research-synthesis/Turing_Degrees_Synthesis.tex}.

\bibitem{RepoReadme}
V.~Reshetnikov, ProveIt repository,
\emph{Turing degrees in Lean and Rocq}, project README, same snapshot.
\url{https://github.com/VladimirReshetnikov/ProveIt/blob/1a1396d4d3a2ac6812692df520517d86ba3a4785/Computability/TuringDegrees/README.md}.

\bibitem{RepoHyper}
V.~Reshetnikov, ProveIt repository,
\emph{Coarse Classes, Hyperdegrees, and the Cone Filter}, research report 10,
README dated 18 September 2026 with additions of 19 September 2026,
same snapshot. Cited only for its stated higher-computability research
boundary; not used as an input to any theorem here.
\url{https://github.com/VladimirReshetnikov/ProveIt/blob/1a1396d4d3a2ac6812692df520517d86ba3a4785/Computability/TuringDegrees/Research/CoarseDegrees/research-reports/10/README.md}.

\end{thebibliography}
\end{document}
